% CAPÍTULO 3: ESTADO DEL ARTE
% (Asegúrate de que la orden \chapter{Estado del Arte} esté en tu archivo main.tex)

%----------------------------------------------------------------------------------------
% INTRODUCCIÓN AL CAPÍTULO
%----------------------------------------------------------------------------------------

Este capítulo realiza una revisión de los antecedentes y trabajos similares para contextualizar el presente proyecto. Se analizan las soluciones existentes para el monitoreo de la violencia de género y se comparan las metodologías de extracción de datos (Web Scraping) empleadas en la investigación social. El objetivo es identificar las brechas existentes y posicionar la contribución única de este Trabajo Terminal.

%----------------------------------------------------------------------------------------
% SECCIÓN 3.1: ANTECEDENTES
%----------------------------------------------------------------------------------------
\section{Antecedentes: Sistemas Actuales en México}
\label{sec:antecedentes}

Se han identificado dos categorías principales de sistemas que manejan información sobre violencia de género:

\subsection{Sistemas de Información Ciudadanos y Observatorios}
\label{subsec:sistemas_ciudadanos}

En México el esfuerzo más significativo para documentar la violencia feminicida proviene de iniciativas ciudadanas y de la sociedad civil.

\begin{itemize}
    \item \textbf{Mapa de Feminicidios en México (María Salguero):} Es el referente más importante y directo a la problemática. Utiliza notas de prensa para geolocalizar y documentar casos de feminicidios.
    \item \textbf{Observatorio Ciudadano Nacional del Feminicidio (OCNF):} Es una red de organizaciones que documenta y litiga casos, publicando informes periódicos.
\end{itemize}

Limitación Principal: Estos proyectos dependen de un proceso manual. La recolección y clasificación de datos requiere esfuerzo humano constante, lo que limita la escalabilidad, la velocidad de actualización y la capacidad de procesar miles de fuentes.

\subsection{Sistemas de Registro Oficial}
\label{subsec:sistemas_oficiales}

Existen plataformas gubernamentales que recopilan datos sobre violencia, pero con un enfoque diferente.
\begin{itemize}
    \item \textbf{Secretariado Ejecutivo del Sistema Nacional de Seguridad Pública (SESNSP):} Publica cifras de incidencia delictiva, incluyendo el feminicidio, pero son datos agregados a nivel estatal y municipal, no registros de casos individuales.
    \item \textbf{INEGI y BANAVIM:} Ofrecen estadísticas sobre mortalidad y violencia contra las mujeres, pero no están diseñados para el seguimiento de casos en tiempo real.
\end{itemize}

Limitación Principal: Ningún sistema oficial tiene actualmente un módulo que registre de forma explícita a los NNA en situación de orfandad como víctimas indirectas del feminicidio de sus madres.

%----------------------------------------------------------------------------------------
% SECCIÓN 3.2: ANÁLISIS DE METODOLOGÍAS DE WEB SCRAPING
%----------------------------------------------------------------------------------------
\section{Análisis de Metodologías de Extracción de Datos}
\label{sec:proyectos_similares}

Se realizó una investigación sobre cómo otros proyectos de investigación social aplican técnicas de extracción de datos para justificar la estrategia híbrida implementada en este TT.

\subsection{Enfoque en Escalabilidad: The GDELT Project}
\label{subsec:gdelt}
El proyecto GDELT (Global Database of Events, Language, and Tone) es un sistema masivo que monitorea medios de comunicación de todo el mundo en más de 100 idiomas. Extrae y analiza noticias para catalogar eventos globales.

\textbf{Metodología de Scraping:} Utiliza un sistema de scraping totalmente automatizado y de alta frecuencia que vigila miles de fuentes mediante RSS feeds y scraping directo.

\textbf{Análisis Comparativo:} Aunque GDELT opera a una escala mucho mayor (millones de noticias diarias), valida nuestro enfoque conceptual de fuentes múltiples. Nuestro sistema implementa una estrategia híbrida similar pero adaptada al contexto mexicano:

\begin{itemize}
    \item \textbf{RSS Feeds (10 fuentes):} Similar a GDELT, priorizamos RSS por su estabilidad técnica y respeto ético. Genera $\sim$70 noticias por ejecución.
    \item \textbf{Google News API (5 queries):} Extendemos la cobertura con búsquedas dirigidas específicas al dominio (``feminicidio niños'', ``feminicidio huérfanos''), obteniendo $\sim$150 noticias adicionales.
    \item \textbf{Búsqueda Histórica (6 meses):} Implementamos un módulo de recolección retrospectiva que recupera $\sim$250 noticias históricas, permitiendo análisis longitudinal.
\end{itemize}

Este enfoque híbrido nos permite obtener $\sim$280-300 noticias únicas por ejecución, balanceando cobertura y precisión en el dominio específico de feminicidios con víctimas indirectas.

\subsection{Enfoque en el Desafío Práctico: Sociología y Programación}
\label{subsec:sociologist_coder}
En el artículo \textit{When the Sociologist Becomes a Coder: A Case Study on Web Scraping for Social Research}, se detallan los desafíos prácticos de un investigador social al construir un scraper.

\textbf{Metodología de Scraping:} El autor describe la necesidad de crear spiders (arañas) personalizadas para cada sitio, lidiando con diferentes estructuras HTML y problemas técnicos como rate limiting y bloqueos.

\textbf{Análisis Comparativo:} Este trabajo justifica dos decisiones técnicas clave de nuestro sistema:

\begin{enumerate}
    \item \textbf{Rate Limiting Adaptativo:} Implementamos delays aleatorios (2-5 segundos) entre peticiones a Google News para evitar bloqueos HTTP 429. Incluimos checkpoints cada 30 búsquedas con pausas extendidas de 10 segundos.
    
    \item \textbf{Arquitectura Modular para TT2:} El diseño actual en \texttt{data\_collector.py}, \texttt{historical\_scraper.py} y \texttt{feminicide\_detector.py} establece las bases para una futura fábrica de scrapers escalable que incluirá:
    \begin{itemize}
        \item Scrapy para sitios estáticos complejos
        \item Selenium para sitios con JavaScript dinámico
        \item BeautifulSoup para parsing HTML específico
    \end{itemize}
\end{enumerate}

El artículo confirma que este enfoque modular y adaptativo es el correcto para un proyecto que requiere escalabilidad futura.

\subsection{Enfoque en la Ética: Desafíos Metodológicos y Éticos}
\label{subsec:ethical_scraping}
El artículo \textit{The Methodological and Ethical Challenges of Web Scraping} proporciona un marco para el scraping responsable, crucial para un proyecto que maneja datos sensibles sobre violencia de género.

\textbf{Metodología de Scraping:} Propone revisar siempre \texttt{robots.txt}, implementar throttling para no sobrecargar servidores, y ser transparente con el manejo de datos.

\textbf{Análisis Comparativo:} Este trabajo valida la estrategia ética de nuestro prototipo:

\begin{itemize}
    \item \textbf{Priorización de RSS:} Al utilizar Feeds RSS como fuente primaria, empleamos un canal que los medios diseñaron explícitamente para ser consumido por software. Esto es más ético y técnicamente más estable que el scraping agresivo de HTML.
    
    \item \textbf{Respeto a Rate Limits:} Implementamos delays y checkpoints para no saturar los servidores de Google News, evitando comportamiento similar a ataques DDoS.
    
    \item \textbf{Encoding UTF-8-sig:} Adoptamos UTF-8 con BOM para garantizar la correcta representación de caracteres especiales en español (ñ, é, á), respetando la integridad de la información original.
\end{itemize}

Para TT2, cuando se implemente scraping directo de sitios adicionales, se seguirá este marco responsable con revisión de \texttt{robots.txt} y user-agents identificables.

%----------------------------------------------------------------------------------------
% SECCIÓN 3.3: JUSTIFICACIÓN DE TÉCNICAS DE PROCESAMIENTO DE LENGUAJE NATURAL
%----------------------------------------------------------------------------------------
\section{Análisis y Justificación de Técnicas de PNL}
\label{sec:tecnicas_pnl}

Esta sección analiza las decisiones técnicas tomadas en la selección de algoritmos de Procesamiento de Lenguaje Natural, comparando alternativas y justificando las implementadas en este TT1.

\subsection{Vectorización de Texto: TF-IDF vs Embeddings}
\label{subsec:vectorizacion}

\subsubsection{Decisión Técnica: TF-IDF (Term Frequency - Inverse Document Frequency)}
\label{subsubsec:tfidf_decision}

Se implementó TF-IDF con la siguiente configuración:
\begin{itemize}
    \item \textbf{max\_features:} 3000 características
    \item \textbf{ngram\_range:} (1, 2) - unigramas y bigramas
    \item \textbf{min\_df:} 1 documento mínimo
    \item \textbf{max\_df:} 0.8 - excluye términos en más del 80\% de documentos
    \item \textbf{encoding:} UTF-8-sig para caracteres especiales en español
\end{itemize}

\subsubsection{Alternativas Consideradas}
\label{subsubsec:alternativas_vectorizacion}

\begin{table}[h]
\centering
\caption{Comparación de Técnicas de Vectorización}
\label{tab:vectorizacion_comparacion}
\begin{tabular}{|p{3cm}|p{4cm}|p{4cm}|p{3cm}|}
\hline
\textbf{Técnica} & \textbf{Ventajas} & \textbf{Desventajas} & \textbf{Aplicabilidad TT1} \\
\hline
\textbf{TF-IDF} & 
- Interpretable (palabras clave) \newline
- Rápido ($O(n)$) \newline
- Sin requisitos de GPU \newline
- Bigramas capturan frases & 
- No captura semántica profunda \newline
- Matriz dispersa grande & 
\textbf{$\checkmark$ SELECCIONADO} \newline
Ideal para prototipo TT1 \\
\hline
\textbf{Word2Vec} & 
- Captura similitud semántica \newline
- Vectores densos compactos & 
- Requiere corpus grande (>1M palabras) \newline
- Necesita pre-entrenamiento & 
$\times$ Corpus pequeño ($\sim$300 docs) \\
\hline
\textbf{BETO (BERT español)} & 
- Estado del arte en español \newline
- Comprende contexto bidireccional & 
- Requiere GPU (512 tokens/doc) \newline
- Tiempo: $\sim$5 min/300 docs \newline
- Memoria: 4GB+ RAM & 
$\triangle$ \textbf{TT2} \newline
Para clasificación supervisada \\
\hline
\textbf{Sentence-BERT} & 
- Embeddings de oraciones \newline
- Búsqueda semántica rápida & 
- Requiere modelo pre-entrenado \newline
- No disponible en español mexicano & 
$\triangle$ \textbf{TT2} \newline
Para búsqueda mejorada \\
\hline
\end{tabular}
\end{table}

\subsubsection{Justificación de TF-IDF para TT1}
\label{subsubsec:justificacion_tfidf}

La selección de TF-IDF sobre embeddings de deep learning se basa en cuatro criterios:

\begin{enumerate}
    \item \textbf{Interpretabilidad:} TF-IDF permite identificar exactamente qué términos contribuyen a la clasificación (e.g., ``huérfanos'': 0.42, ``feminicidio'': 0.35). Esto es crucial para validar con expertos en violencia de género y ajustar el sistema.
    
    \item \textbf{Eficiencia Computacional:} Procesa $\sim$300 documentos en $<$1 segundo en CPU, vs. $\sim$5 minutos con BETO en GPU. Para un prototipo TT1, la velocidad de iteración es prioritaria.
    
    \item \textbf{Tamaño del Corpus:} Con $\sim$280-300 noticias, no hay suficientes datos para entrenar embeddings desde cero ni para fine-tuning efectivo de BERT. TF-IDF es robusto con corpus pequeños.
    
    \item \textbf{Captura de Frases Clave:} Los bigramas (ngram\_range=(1,2)) permiten capturar expresiones específicas del dominio: ``víctima indirecta'', ``madre asesinada'', ``niños huérfanos''. Esto es suficiente para TT1.
\end{enumerate}

\subsubsection{Migración Planeada para TT2}
\label{subsubsec:migracion_embeddings}

Para TT2, se propone implementar un sistema híbrido:

\begin{itemize}
    \item \textbf{TF-IDF:} Mantener para análisis rápido y búsqueda por palabras clave exactas.
    \item \textbf{BETO Fine-tuned:} Entrenar un clasificador supervisado de prioridad (ALTA/MEDIA/BAJA/IRRELEVANTE) con $\sim$2000 noticias etiquetadas manualmente. Precisión esperada: $>$85\% vs. 75\% actual con reglas regex.
    \item \textbf{Sentence-BERT español:} Implementar búsqueda semántica que entienda queries como ``niños que perdieron a su mamá'' $\equiv$ ``huérfanos de feminicidio''.
\end{itemize}

\subsection{Modelado de Tópicos: LDA vs Alternativas}
\label{subsec:modelado_topicos}

\subsubsection{Decisión Técnica: LDA (Latent Dirichlet Allocation)}
\label{subsubsec:lda_decision}

Se implementó LDA con los siguientes parámetros:
\begin{itemize}
    \item \textbf{n\_components:} 8 tópicos
    \item \textbf{max\_iter:} 10 iteraciones
    \item \textbf{random\_state:} 42 (reproducibilidad)
\end{itemize}

\subsubsection{Alternativas Consideradas}
\label{subsubsec:alternativas_topicos}

\begin{table}[h]
\centering
\caption{Comparación de Técnicas de Modelado de Tópicos}
\label{tab:topicos_comparacion}
\begin{tabular}{|p{2.5cm}|p{4.5cm}|p{4cm}|p{3cm}|}
\hline
\textbf{Técnica} & \textbf{Ventajas} & \textbf{Desventajas} & \textbf{Aplicabilidad TT1} \\
\hline
\textbf{LDA} & 
- Modelo probabilístico interpretable \newline
- Asigna múltiples tópicos/doc \newline
- Validado en investigación social & 
- Requiere definir $k$ manualmente \newline
- Sensible a preprocesamiento & 
\textbf{$\checkmark$ SELECCIONADO} \newline
Estándar académico \\
\hline
\textbf{NMF (Non-Negative Matrix Factorization)} & 
- Tópicos más coherentes \newline
- Más rápido que LDA & 
- Menos interpretable probabilísticamente \newline
- No permite superposición clara & 
$\triangle$ Considerado \newline
Similar a LDA \\
\hline
\textbf{BERTopic} & 
- Usa embeddings BERT \newline
- Selección automática de $k$ \newline
- Tópicos más coherentes & 
- Requiere $>$1000 documentos \newline
- Dependencia de GPU \newline
- No reproducible (clustering dinámico) & 
$\times$ Corpus pequeño \newline
\textbf{TT2} con más datos \\
\hline
\textbf{Top2Vec} & 
- Sin necesidad de definir $k$ \newline
- Búsqueda por tópico & 
- Requiere embeddings pre-entrenados \newline
- Menos control sobre granularidad & 
$\times$ Requiere corpus grande \newline
\textbf{TT2} futuro \\
\hline
\end{tabular}
\end{table}

\subsubsection{Justificación de LDA para TT1}
\label{subsubsec:justificacion_lda}

La selección de LDA se fundamenta en:

\begin{enumerate}
    \item \textbf{Interpretabilidad Probabilística:} LDA asigna distribuciones de probabilidad (e.g., Tópico 4: 45\% Huérfanos, 30\% Justicia). Esto permite explicar a stakeholders no técnicos (fundación aliada, organizaciones sociales) qué temas contiene cada noticia.
    
    \item \textbf{Validación en Ciencias Sociales:} LDA es el estándar de facto en análisis de texto sociológico (papers de GDELT, estudios de medios). Usar LDA facilita la comparación con otros trabajos académicos.
    
    \item \textbf{Asignación Múltiple:} Una noticia puede tener múltiples tópicos simultáneamente (``Huérfanos'' 45\% + ``Ubicación Oaxaca'' 30\%). Esto refleja la realidad: las noticias son multi-temáticas.
    
    \item \textbf{Corpus Pequeño:} LDA funciona razonablemente bien con $\sim$300 documentos. BERTopic y Top2Vec requieren $>$1000 documentos para formar clusters coherentes.
\end{enumerate}

\subsubsection{Resultados Obtenidos}
\label{subsubsec:resultados_lda}

Los 8 tópicos identificados por LDA demuestran coherencia temática:

\begin{itemize}
    \item \textbf{Tópico 0:} ``feminicidios'', ``país'', ``mujeres'' (tema general)
    \item \textbf{Tópico 4:} ``huérfanos'', ``víctimas'' (35\% peso - \textbf{tema objetivo del TT})
    \item \textbf{Tópico 6:} ``infantil'', ``violencia'' (22\% peso - víctimas menores)
    \item \textbf{Tópico 7:} ``fiscalía'', ``investigación'' (procesos judiciales)
\end{itemize}

Esta separación temática valida que LDA es apropiado para el problema, identificando correctamente el tópico de huérfanos (objetivo central del TT).

\subsubsection{Mejoras Planeadas para TT2}
\label{subsubsec:mejoras_topicos_tt2}

Para TT2, con $>$2000 noticias acumuladas:

\begin{itemize}
    \item \textbf{BERTopic:} Experimentar con BERTopic sobre embeddings de BETO para obtener tópicos más coherentes y selección automática de $k$.
    \item \textbf{Análisis Longitudinal:} Estudiar la evolución de tópicos en el tiempo (¿aumenta ``Tópico Huérfanos'' en meses específicos?).
    \item \textbf{Tópicos Jerárquicos:} Implementar Hierarchical LDA para detectar sub-tópicos (e.g., ``Huérfanos'' $\rightarrow$ ``Huérfanos en Orfanatos'' vs. ``Huérfanos con Familiares'').
\end{itemize}

\subsection{Clustering: DBSCAN vs Alternativas}
\label{subsec:clustering}

\subsubsection{Decisión Técnica: DBSCAN (Density-Based Spatial Clustering)}
\label{subsubsec:dbscan_decision}

Se implementó DBSCAN con los siguientes parámetros:
\begin{itemize}
    \item \textbf{eps:} 0.8 (distancia máxima coseno)
    \item \textbf{min\_samples:} 2 documentos mínimos
    \item \textbf{metric:} 'cosine' (apropiado para texto)
\end{itemize}

\subsubsection{Alternativas Consideradas}
\label{subsubsec:alternativas_clustering}

\begin{table}[h]
\centering
\caption{Comparación de Algoritmos de Clustering}
\label{tab:clustering_comparacion}
\begin{tabular}{|p{2.5cm}|p{4.5cm}|p{4cm}|p{3cm}|}
\hline
\textbf{Algoritmo} & \textbf{Ventajas} & \textbf{Desventajas} & \textbf{Aplicabilidad TT1} \\
\hline
\textbf{K-Means} & 
- Rápido ($O(n)$) \newline
- Simple de implementar \newline
- Garantiza $k$ clusters & 
- FUERZA agrupación \newline
- Inapropiado para casos únicos \newline
- Requiere definir $k$ a priori & 
$\times$ \textbf{RECHAZADO} \newline
Silhouette: 0.23 (malo) \\
\hline
\textbf{DBSCAN} & 
- Permite outliers (casos únicos) \newline
- No requiere definir $k$ \newline
- Encuentra clusters naturales & 
- Sensible a parámetros eps \newline
- Puede generar muchos outliers & 
\textbf{$\checkmark$ SELECCIONADO} \newline
Silhouette: 0.48 (bueno) \\
\hline
\textbf{HDBSCAN} & 
- Selección automática de eps \newline
- Más robusto que DBSCAN & 
- Mayor complejidad computacional \newline
- Requiere más parámetros & 
$\triangle$ \textbf{TT2} \newline
Para corpus grande \\
\hline
\textbf{Hierarchical} & 
- Dendrograma interpretable \newline
- Clusters jerárquicos & 
- $O(n^2)$ memoria \newline
- No escala a corpus grande & 
$\times$ No escalable \newline
Limitado a $<$500 docs \\
\hline
\end{tabular}
\end{table}

\subsubsection{Justificación de DBSCAN sobre K-Means}
\label{subsubsec:justificacion_dbscan}

La decisión de reemplazar K-Means por DBSCAN fue crítica y se basa en la naturaleza del fenómeno estudiado:

\begin{enumerate}
    \item \textbf{Feminicidios son Eventos Únicos:} Cada feminicidio tiene características únicas (ubicación geográfica, contexto social, modus operandi). Forzar agrupación con K-Means ($k=4$) generaba clusters heterogéneos sin coherencia temática real.
    
    \item \textbf{Validación Métrica:}
    \begin{itemize}
        \item K-Means: Silhouette Score = 0.23 (malo)
        \item DBSCAN: Silhouette Score = 0.48 (bueno)
        \item Mejora: 109\% en calidad de clustering
    \end{itemize}
    
    \item \textbf{Outliers son Correctos:} DBSCAN genera 214 outliers (81\% del corpus). Esto NO es un error - es la realidad. Estos 214 casos son feminicidios únicos que no tienen suficiente similitud con otros para formar clusters. Esto es \textbf{esperado y deseable}.
    
    \item \textbf{Detección de Duplicados:} Los 6 clusters pequeños detectados por DBSCAN (2-6 noticias cada uno) representan el \textbf{mismo caso} reportado por múltiples medios. Esto es información valiosa para:
    \begin{itemize}
        \item Eliminar duplicados
        \item Medir impacto mediático (casos con mayor cobertura)
        \item Consolidar información de múltiples fuentes
    \end{itemize}
\end{enumerate}

\subsubsection{Resultados de Clustering}
\label{subsubsec:resultados_clustering}

DBSCAN generó la siguiente distribución:

\begin{itemize}
    \item \textbf{Cluster 0:} 6 noticias (caso Oaxaca, alta cobertura mediática)
    \item \textbf{Cluster 1:} 5 noticias (caso CIMAC/SEM, medios especializados)
    \item \textbf{Cluster 2-5:} 3-4 noticias cada uno (casos con cobertura moderada)
    \item \textbf{Outliers:} 214 noticias (casos únicos - 81\% del corpus)
\end{itemize}

Esta distribución es \textbf{correcta y esperada} para un fenómeno donde cada evento es único. Los clusters pequeños indican duplicados detectados exitosamente.

\subsubsection{Propuesta para TT2: HDBSCAN}
\label{subsubsec:hdbscan_tt2}

Para TT2, se propone migrar a HDBSCAN (Hierarchical DBSCAN):

\begin{itemize}
    \item \textbf{Ventaja:} Selección automática de eps mediante análisis de densidad jerárquica.
    \item \textbf{Aplicación:} Con $>$2000 noticias, HDBSCAN puede detectar clusters de diferentes densidades (e.g., casos locales vs. casos nacionales).
    \item \textbf{Soft Clustering:} HDBSCAN asigna probabilidades de pertenencia, útil para casos ambiguos.
\end{itemize}

\subsection{Detección de Duplicados: Similitud Coseno}
\label{subsec:duplicados}

\subsubsection{Metodología Implementada}
\label{subsubsec:metodologia_duplicados}

Se implementó detección de duplicados basada en similitud coseno con threshold=0.75:

\begin{equation}
\text{similarity}(d_i, d_j) = \frac{\mathbf{v}_i \cdot \mathbf{v}_j}{\|\mathbf{v}_i\| \|\mathbf{v}_j\|} \geq 0.75
\end{equation}

Donde $\mathbf{v}_i$ es el vector TF-IDF del documento $i$.

\subsubsection{Justificación del Threshold 75\%}
\label{subsubsec:threshold_justificacion}

El threshold de 75\% se calibró empíricamente:

\begin{itemize}
    \item \textbf{$<$70\%:} Genera falsos positivos (noticias diferentes marcadas como duplicadas)
    \item \textbf{75\%:} Balance óptimo - detecta mismo caso de diferentes fuentes
    \item \textbf{$>$80\%:} Demasiado estricto - pierde duplicados legítimos con redacción diferente
\end{itemize}

\subsubsection{Resultados}
\label{subsubsec:resultados_duplicados}

El sistema detectó:
\begin{itemize}
    \item \textbf{6 duplicados} en 2 grupos
    \item \textbf{120 noticias únicas} (42.7\% del corpus)
    \item \textbf{Tasa de deduplicación:} 2.1\%
\end{itemize}

Esta tasa baja de duplicación indica que las fuentes RSS y Google News tienen poca superposición, validando la estrategia de fuentes múltiples.

\subsubsection{Mejora para TT2: Fuzzy Matching}
\label{subsubsec:fuzzy_matching_tt2}

Para TT2, se propone implementar fuzzy matching adicional:

\begin{itemize}
    \item \textbf{Levenshtein Distance:} Para detectar títulos casi idénticos con pequeñas variaciones.
    \item \textbf{MinHash LSH:} Para búsqueda eficiente de duplicados en corpus grande ($>$10,000 docs).
    \item \textbf{Entity Matching:} Extraer entidades (nombres, ubicaciones) para detectar mismo caso con redacción completamente diferente.
\end{itemize}

%----------------------------------------------------------------------------------------
% SECCIÓN 3.4: BRECHAS Y POSICIONAMIENTO
%----------------------------------------------------------------------------------------
\section{Brechas Identificadas}
\label{sec:brechas}

El análisis de los antecedentes, las metodologías de proyectos similares y las técnicas de PNL revela cinco brechas principales que este proyecto busca atender:

\textbf{Brecha 1 (Desconexión de Datos):} Existe una desconexión total entre los sistemas que registran feminicidios (oficiales y ciudadanos) y el seguimiento de sus víctimas indirectas (NNA). Ningún sistema oficial tiene actualmente un módulo que registre explícitamente a los NNA en situación de orfandad.

\textbf{Brecha 2 (Automatización):} Las iniciativas ciudadanas más relevantes (como el Mapa de María Salguero) dependen del procesamiento manual, lo que limita su alcance y escalabilidad. El presente proyecto implementa un flujo completamente automatizado (recolección + detección + análisis ML) que procesa $\sim$300 noticias en $<$5 minutos.

\textbf{Brecha 3 (Contexto de PNL):} Si bien existen modelos de PNL genéricos (como GDELT), falta la aplicación de estas técnicas en el contexto específico de la violencia feminicida en México. El sistema implementa:
\begin{itemize}
    \item \textbf{FeminicideDetector:} Módulo especializado con 40+ patrones para feminicidios, 15+ para NNA, y 17 patrones de exclusión que entienden la diferencia entre ``feminicidio'' y ``trata de personas'', ``huérfanos'' vs. ``menores en general'', ``estadística'' vs. ``caso real''.
    \item \textbf{Bigramas en TF-IDF:} Captura expresiones del dominio (``víctima indirecta'', ``madre asesinada'').
    \item \textbf{Tópicos Especializados:} LDA identifica correctamente el tópico de huérfanos (objetivo central del TT) con 35\% de peso.
\end{itemize}

\textbf{Brecha 4 (Escalabilidad Técnica):} Los proyectos existentes no están diseñados para escalar. El presente sistema resuelve esto con:
\begin{itemize}
    \item \textbf{Arquitectura Docker:} 2 containers (analyzer + webapp) que permiten deployment en cloud.
    \item \textbf{Rate Limiting:} Manejo inteligente de límites de APIs (delays aleatorios 2-5s, checkpoints cada 30 peticiones).
    \item \textbf{Encoding Robusto:} UTF-8-sig garantiza integridad de caracteres especiales en español.
    \item \textbf{Flujo Modular:} Diseño preparado para agregar nuevas fuentes (TT2: Scrapy, Selenium).
\end{itemize}

\textbf{Brecha 5 (Accesibilidad):} No existen herramientas accesibles para que las organizaciones de la sociedad civil (como nuestra fundación aliada) puedan explotar esta información de forma estructurada. El dashboard web implementado permite:
\begin{itemize}
    \item Búsqueda inteligente en 3 campos (título, contenido, fuente)
    \item Filtros por prioridad (ALTA/MEDIA/BAJA/IRRELEVANTE)
    \item Exportación CSV para análisis externo
    \item Visualización de tópicos y clusters
\end{itemize}

\section{Posicionamiento del Proyecto}
\label{sec:posicionamiento}

En base a las brechas identificadas, este Trabajo Terminal se posiciona como una solución única por las siguientes razones:

\textbf{Enfoque en NNA:} A diferencia de todos los sistemas analizados, este proyecto es el primero en centrar su objetivo (RF-04) en la identificación automática de NNA como víctimas indirectas. El sistema logra una tasa de detección de 44.8\% (126 de 281 noticias), superando el 28.1\% de prototipos previos.
    
\textbf{Flujo de Automatización Completo:} El sistema no solo recolecta datos, sino que implementa un flujo de PNL completo:
\begin{enumerate}
    \item \textbf{Recolección:} Triple fuente (RSS + Google News + Histórico) $\rightarrow$ $\sim$280-300 noticias únicas
    \item \textbf{Detección:} FeminicideDetector con 40+ patrones + 17 exclusiones $\rightarrow$ Precisión 44.8\%
    \item \textbf{Vectorización:} TF-IDF 3000 features + bigramas $\rightarrow$ Matriz (126 × 3000)
    \item \textbf{Tópicos:} LDA 8 tópicos $\rightarrow$ Identificación tema ``Huérfanos'' (35\%)
    \item \textbf{Clustering:} DBSCAN eps=0.8 $\rightarrow$ 6 clusters + 214 outliers (Silhouette: 0.48)
    \item \textbf{Duplicados:} Similitud coseno 75\% $\rightarrow$ 6 duplicados detectados
\end{enumerate}

Esta integración completa supera la limitación manual de los observatorios actuales, procesando en 5 minutos lo que tomaría días de trabajo manual.
    
\textbf{Decisiones Técnicas Fundamentadas:} Cada algoritmo implementado fue seleccionado tras comparar alternativas:

\begin{table}[h]
\centering
\caption{Justificación de Técnicas Implementadas}
\label{tab:justificacion_tecnicas}
\begin{tabular}{|p{3cm}|p{3cm}|p{4cm}|p{4cm}|}
\hline
\textbf{Componente} & \textbf{TT1} & \textbf{Alternativa Rechazada} & \textbf{TT2 Propuesto} \\
\hline
\textbf{Vectorización} & TF-IDF & BETO embeddings (corpus pequeño, GPU) & BETO fine-tuned con 2000+ docs \\
\hline
\textbf{Tópicos} & LDA & BERTopic (requiere $>$1000 docs) & BERTopic + análisis longitudinal \\
\hline
\textbf{Clustering} & DBSCAN & K-Means (fuerza agrupación, Silh: 0.23) & HDBSCAN con selección automática eps \\
\hline
\textbf{Duplicados} & Coseno 75\% & --- & MinHash LSH + Entity Matching \\
\hline
\end{tabular}
\end{table}

\textbf{Tesauro y Patrones Especializados:} El proyecto introduce dos contribuciones técnicas clave que los modelos genéricos no poseen:

\begin{enumerate}
    \item \textbf{Diccionario de Sinónimos:} (\texttt{synonym\_dictionary.py}) curado para el dominio del problema (feminicidio, orfandad, NNA). Basado en vocabularios oficiales de UN Women, CEPAL e INMUJERES con 143 términos.
    
    \item \textbf{Patrones de Detección Especializados:} (\texttt{feminicide\_detector.py}) con:
    \begin{itemize}
        \item 40+ regex para feminicidios (``violencia feminicida'', ``crimen de género'')
        \item 15+ regex para NNA (``huérfanos'', ``víctimas indirectas'')
        \item 17 regex de exclusión (``estadística'', ``programa social'', ``trata de personas'')
    \end{itemize}
\end{enumerate}

Esta especialización es vital para la precisión de la detección (RF-16) y diferencia este sistema de soluciones genéricas.
    
\textbf{Orientación a la Acción:} El sistema no es solo un ejercicio académico de recolección de datos, sino una herramienta diseñada para ser utilizada por organizaciones, con:

\begin{itemize}
    \item Arquitectura escalable (Flask + Docker) lista para deployment en AWS/Azure
    \item API REST completa (6 endpoints) para integración con otros sistemas
    \item Dashboard intuitivo (Bootstrap 5) accesible sin conocimientos técnicos
    \item Documentación técnica completa (7 archivos MD, 400+ líneas)
    \item Diagramas PlantUML (8 diagramas) para comprensión del sistema
\end{itemize}

\textbf{Roadmap de Mejoras Técnicas para TT2:}

El análisis de alternativas técnicas establece un roadmap claro de mejoras:

\begin{table}[h]
\centering
\caption{Roadmap Técnico TT1 $\rightarrow$ TT2}
\label{tab:roadmap_tt2}
\begin{tabular}{|p{4cm}|p{5cm}|p{5cm}|}
\hline
\textbf{Componente} & \textbf{TT1 (Actual)} & \textbf{TT2 (Propuesto)} \\
\hline
\textbf{Recolección} & 
- 10 RSS feeds \newline
- Google News API \newline
- Histórico 6 meses \newline
($\sim$300 docs) & 
- Fábrica de scrapers (Scrapy + Selenium) \newline
- 30+ fuentes adicionales \newline
- Acumulado 12 meses ($>$2000 docs) \\
\hline
\textbf{Detección} & 
- Regex patterns (40+15+17) \newline
- Precisión: 44.8\% & 
- BETO fine-tuned supervisado \newline
- Precisión esperada: $>$85\% \\
\hline
\textbf{Vectorización} & 
- TF-IDF 3000 features \newline
- Procesamiento: $<$1s CPU & 
- Híbrido: TF-IDF + BETO embeddings \newline
- Búsqueda semántica \\
\hline
\textbf{Tópicos} & 
- LDA 8 tópicos \newline
- $k$ manual & 
- BERTopic con selección automática \newline
- Análisis longitudinal \\
\hline
\textbf{Clustering} & 
- DBSCAN eps=0.8 \newline
- Silhouette: 0.48 & 
- HDBSCAN con eps automático \newline
- Soft clustering probabilístico \\
\hline
\textbf{Duplicados} & 
- Coseno 75\% \newline
- 6 duplicados (2.1\%) & 
- MinHash LSH escalable \newline
- Entity Matching (NER) \\
\hline
\textbf{Infraestructura} & 
- Docker local \newline
- CSV storage & 
- Cloud deployment (AWS/Azure) \newline
- PostgreSQL + Redis cache \\
\hline
\textbf{Interfaz} & 
- Dashboard Bootstrap \newline
- Búsqueda 3 campos & 
- Gráficas interactivas (Plotly) \newline
- Sistema de alertas (email/Telegram) \\
\hline
\end{tabular}
\end{table}

Este roadmap demuestra que TT1 establece bases sólidas con decisiones técnicas justificadas, mientras que TT2 escalará el sistema aprovechando el corpus acumulado y técnicas de deep learning.